%==============================================================================
% Chapitre 23 - Introduction aux projectiles
%==============================================================================

\section{Introduction}

Le projectile constitue l'élément mobile du système d'arme.

Son rôle est de transporter une énergie ou un effet jusqu'à une cible.

Bien que souvent désigné sous le terme générique de « balle », le projectile
peut prendre de nombreuses formes et remplir des fonctions très différentes.

L'étude des projectiles est au cœur de la balistique.

Leur géométrie, leur masse, leur matériau et leur stabilité influencent
directement les performances observées en vol.

\begin{aretenir}

Le projectile est l'objet mis en mouvement par l'arme et étudié par la
balistique.

\end{aretenir}

\section{Définition}

Un projectile est un corps destiné à être propulsé vers une cible.

Il peut être lancé par différents moyens :

\begin{itemize}
    \item énergie chimique ;
    \item énergie pneumatique ;
    \item énergie mécanique ;
    \item énergie électromagnétique.
\end{itemize}

Dans ce manuel, nous nous intéresserons principalement aux projectiles
propulsés par des armes à feu.

\section{Projectile et munition}

Les termes projectile et munition sont souvent confondus.

Ils désignent pourtant des éléments différents.

La munition comprend généralement :

\begin{itemize}
    \item le projectile ;
    \item l'étui ;
    \item la charge propulsive ;
    \item l'amorce.
\end{itemize}

Le projectile n'est donc qu'un composant de la munition.

\begin{aretenir}

Une balle n'est pas une cartouche.

Le projectile n'est qu'un des éléments constituant la munition.

\end{aretenir}

\section{Fonctions du projectile}

Selon sa conception, un projectile peut avoir plusieurs fonctions :

\begin{itemize}
    \item transmettre de l'énergie cinétique ;
    \item perforer une cible ;
    \item produire un effet incendiaire ;
    \item produire un effet traceur ;
    \item transporter une charge utile ;
    \item produire un effet explosif.
\end{itemize}

Certaines munitions combinent plusieurs fonctions.

\section{Cycle de vie d'un projectile}

Du point de vue balistique, la vie d'un projectile peut être divisée en
plusieurs phases :

\begin{enumerate}
    \item stockage ;
    \item mise à feu ;
    \item déplacement dans le canon ;
    \item sortie de bouche ;
    \item vol ;
    \item impact.
\end{enumerate}

Ces phases correspondent aux différentes branches de la balistique :

\begin{itemize}
    \item balistique intérieure ;
    \item balistique intermédiaire ;
    \item balistique extérieure ;
    \item balistique terminale.
\end{itemize}

\section{Évolution historique}

Les premiers projectiles étaient généralement sphériques.

Ils étaient souvent fabriqués en :

\begin{itemize}
    \item pierre ;
    \item plomb ;
    \item fonte.
\end{itemize}

L'apparition des canons rayés et l'augmentation des vitesses ont conduit au
développement de formes plus aérodynamiques.

Les projectiles modernes présentent généralement :

\begin{itemize}
    \item une forme allongée ;
    \item une meilleure stabilité ;
    \item une traînée réduite.
\end{itemize}

\section{Pourquoi les projectiles ne sont-ils plus sphériques ?}

Une sphère présente plusieurs inconvénients :

\begin{itemize}
    \item forte traînée ;
    \item stabilité limitée ;
    \item mauvaise conservation de vitesse.
\end{itemize}

Les projectiles modernes utilisent des formes permettant :

\begin{itemize}
    \item une meilleure pénétration dans l'air ;
    \item une meilleure stabilité ;
    \item une portée accrue.
\end{itemize}

\section{Les grandes caractéristiques}

Plusieurs paramètres permettent de décrire un projectile :

\begin{itemize}
    \item calibre ;
    \item masse ;
    \item longueur ;
    \item géométrie ;
    \item matériau ;
    \item vitesse ;
    \item énergie.
\end{itemize}

Ces paramètres influencent directement le comportement en vol.

\section{Projectile et énergie}

Le projectile transporte une partie de l'énergie produite lors du tir.

Cette énergie est principalement de nature cinétique.

\begin{equation}
E_c = \frac{1}{2}mv^2
\end{equation}

où :

\begin{itemize}
    \item $E_c$ est l'énergie cinétique ;
    \item $m$ est la masse du projectile ;
    \item $v$ est sa vitesse.
\end{itemize}

\begin{aretenir}

La vitesse influence l'énergie au carré.

Une faible variation de vitesse peut donc produire une variation importante
de l'énergie.

\end{aretenir}

\section{Projectile et aérodynamique}

Pendant son vol, le projectile interagit en permanence avec l'atmosphère.

Il est soumis notamment :

\begin{itemize}
    \item à la gravité ;
    \item à la traînée ;
    \item au vent ;
    \item à divers effets aérodynamiques.
\end{itemize}

Son comportement dépend fortement de sa géométrie.

\section{Projectile et stabilité}

Un projectile doit conserver une orientation compatible avec son vol.

La stabilité constitue donc une caractéristique fondamentale.

Elle dépend notamment :

\begin{itemize}
    \item de la géométrie ;
    \item de la masse ;
    \item de la rotation ;
    \item des forces aérodynamiques.
\end{itemize}

Cette notion sera étudiée en détail dans les chapitres suivants.

\section{Classification générale}

Les projectiles peuvent être classés selon différents critères.

\subsection{Selon la fonction}

\begin{itemize}
    \item ordinaires ;
    \item perforants ;
    \item incendiaires ;
    \item traceurs ;
    \item explosifs.
\end{itemize}

\subsection{Selon la forme}

\begin{itemize}
    \item sphériques ;
    \item cylindro-ogivaux ;
    \item flèches ;
    \item formes spécialisées.
\end{itemize}

\subsection{Selon la stabilisation}

\begin{itemize}
    \item stabilisés par rotation ;
    \item stabilisés par empennage ;
    \item stabilisation mixte.
\end{itemize}

\section{Projectiles et simulation}

Les modèles balistiques représentent le projectile à différents niveaux de
détail.

Selon les besoins, il peut être décrit par :

\begin{itemize}
    \item une simple masse ponctuelle ;
    \item un corps rigide ;
    \item un modèle aérodynamique complet.
\end{itemize}

Le choix dépend du niveau de fidélité recherché.

\begin{simulation}

Les modèles les plus simples ignorent totalement la forme du projectile.

Les modèles avancés utilisent sa géométrie et ses caractéristiques
aérodynamiques pour calculer son comportement en vol.

\end{simulation}

\section{Vers les chapitres suivants}

Les prochains chapitres étudieront les propriétés fondamentales des
projectiles :

\begin{itemize}
    \item géométrie ;
    \item masse ;
    \item matériaux ;
    \item centres caractéristiques ;
    \item stabilité ;
    \item aérodynamique.
\end{itemize}

Ces notions permettront ensuite de comprendre les phénomènes plus complexes
comme la dérive gyroscopique ou le coefficient balistique.

\section{À retenir}

\begin{aretenir}

\begin{itemize}
    \item Le projectile est l'objet propulsé vers la cible.
    \item Il ne doit pas être confondu avec la munition.
    \item Sa géométrie influence fortement ses performances.
    \item Il transporte principalement de l'énergie cinétique.
    \item Sa stabilité constitue une condition essentielle d'un vol correct.
\end{itemize}

\end{aretenir}
